\section{Introduction}
Social media platforms have become an indispensable part of modern life, profoundly reshaping how individuals disseminate, exchange, and form opinions~\cite{anderson2019recent,dong2018survey,jia2015opinion}. Compared to traditional channels of information diffusion, this transformation has significantly increased the complexity of opinion dynamics and exerted far-reaching effects on human behavior and social interaction patterns~\cite{Le20}. The pervasive nature of online interactions has spurred extensive academic interest in understanding the mechanisms underlying opinion propagation and formation, leading to the development of various mathematical models to capture these processes~\cite{jia2015opinion, PrTe17, dong2018survey, anderson2019recent}. Among them, the Friedkin–Johnsen (FJ) model ~\cite{FrJo90} has emerged as one of the most widely applied frameworks. Notably, adaptations of the FJ model have been used to analyze complex negotiation processes, such as those leading to the Paris Agreement~\cite{BeWaVaHoShAl21}, demonstrating its utility in explaining consensus formation in multilateral international settings.

The core of the FJ model lies in its equilibrium opinion vector, which represents the stable set of expressed opinions reached after prolonged social influence. This equilibrium is determined jointly by the network topology, individuals’ intrinsic opinions, and interpersonal influence weights. It carries clear sociological meaning and serves as the foundation for quantifying key societal phenomena such as polarization~\cite{matakos2017measuring,dandekar2013biased, musco2018minimizing}, disagreement~\cite{dandekar2013biased, musco2018minimizing}, conflict~\cite{chen2018quantifying}, and controversy~\cite{chen2018quantifying}. However, existing solvers~\cite{xu2021fast,neumann2024sublinear,wang2025efficient} typically implement opinion updates as a synchronous weighted average of initial opinions and neighbors’ opinions from the previous iteration. This approach implicitly assumes that individuals adjust their views based solely on outdated information from all neighbors, with no access to opinions updated earlier in the current round and no memory of their own past states. While this simplification facilitates theoretical analysis, it contradicts the temporal dependencies observed in real social interactions, where people often respond immediately to the latest statements and exhibit cognitive continuity over time. More critically, this static update scheme leads to extremely slow convergence on large-scale networks, severely limiting the practical applicability of the FJ model at real-world scales. Although numerous acceleration techniques exist for consensus-only models like DeGroot, they cannot be directly applied here. The FJ model explicitly preserves fixed intrinsic opinions, resulting in a more complex equilibrium structure that defies straightforward adaptation of existing methods.

To address these challenges, we propose a new class of efficient solvers that strictly preserve the original equilibrium of the FJ model while dramatically accelerating convergence through explicit modeling of temporal dynamics. Specifically, our approach captures responsiveness to the latest expressed opinions and retention of historical memory. It integrates three key mechanisms: sequential updates that immediately propagate freshly computed opinions within each iteration, self-interaction effects that enhance stability and enable faster convergence, and second-order historical dependence that incorporates past states to better capture the inertia of opinion evolution. We theoretically prove that all proposed algorithms converge to the same unique equilibrium as the classical model, yet with provably faster convergence rates. Experimental results demonstrate that our proposed algorithms exhibit superior performance and scalability on networks comprising tens of millions of nodes. The main contributions of this paper are summarized as follows:

\wgy{ change this
\begin{itemize}
    \item We introduce a robust local iteration algorithm that efficiently approximates the equilibrium opinion vector while guaranteeing relative error bounds.
    \item By incorporating SOR techniques, we optimize the algorithm with a relaxation parameter $\omega$, leading to faster convergence and improved computational efficiency.
    \item We conduct extensive numerical experiments on real-world network datasets. The results validate the advantages of our algorithms in terms of high efficiency and strong scalability.
\end{itemize}
% The remainder of this paper is organized as follows: Section~\ref{sec:related} reviews related work on opinion dynamics and the FJ model; Section~\ref{sec:prelim} introduces notation and preliminaries necessary for understanding the subsequent sections; Section~\ref{sec:problem} provides an overview of existing algorithms and their limitations; Section~\ref{sec:localIter} details the local iteration algorithm and SOR optimization; Section~\ref{sec:experim} presents experimental results and performance comparisons; and Section~\ref{sec:conclusion} concludes the paper and summarizes the key findings.
}